\documentclass[12pt,a4paper]{article}

\usepackage[utf8]{inputenc}
\usepackage[T1]{fontenc}
\usepackage[brazil]{babel}
\usepackage{lmodern}
\usepackage{microtype}
\usepackage{geometry}
\usepackage{amsmath,amssymb}
\usepackage{booktabs}
\usepackage{longtable}
\usepackage{array}
\usepackage{tabularx}
\usepackage{pdflscape}
\usepackage{url}
\usepackage{hyperref}

\geometry{margin=2.5cm}
\hypersetup{
  colorlinks=true,
  linkcolor=blue,
  citecolor=blue,
  urlcolor=blue
}

\newcolumntype{Y}{>{\raggedright\arraybackslash}X}
\setlength{\parindent}{1.25cm}
\setlength{\parskip}{0.2em}
\setlength{\LTpre}{0.5em}
\setlength{\LTpost}{0.5em}

\title{Otimização Contínua de Pesos de Alocação para Network Slicing em Redes 5G/O-RAN usando Meta-Heurísticas}
\author{Autor(es)}
\date{}

\begin{document}

\maketitle

\begin{abstract}
Este artigo investiga o uso de meta-heurísticas para otimizar a distribuição de recursos entre slices em redes 5G/O-RAN. O problema é formulado como uma otimização contínua, em que uma solução candidata é um vetor de pesos de alocação para os slices eMBB, URLLC e MTC. A avaliação experimental utiliza simulações em ns-3/ns-O-RAN Gym e compara Particle Swarm Optimization (PSO), Genetic Algorithm (GA), Simulated Annealing (SA) e uma meta-heurística híbrida que combina exploração populacional, operadores genéticos, atualização inspirada em PSO e refinamento local estocástico. Os resultados da campanha \texttt{20260623\_123740} mostram que o PSO obteve a melhor solução nos cenários normal e congestionado, com scores de 68,085 e 84,265, respectivamente. A meta-heurística híbrida obteve a melhor solução no cenário de baixo tráfego, com score de 56,818. Considerando os três cenários comparáveis, o PSO foi o método mais consistente em termos de melhor solução encontrada, enquanto a híbrida apresentou desempenho competitivo, mas não superior de forma geral. A comparação pareada com heurísticas de alocação na semente 1 indica ganhos pequenos em baixo tráfego, moderados no cenário normal e relevantes em congestionamento. Entretanto, a evidência estatística ainda é limitada, pois a busca meta-heurística disponível não contém múltiplas execuções independentes por método e os baselines puros RR, BCQI, PF, AQPS e RSLAQ não estão completos no mesmo protocolo experimental da campanha analisada.
\end{abstract}

\noindent\textbf{Palavras-chave:} network slicing; 5G; O-RAN; meta-heurísticas; PSO; algoritmo genético; simulated annealing; alocação de recursos.

\section{Introdução}

Redes 5G introduzem suporte nativo a network slicing, permitindo que diferentes classes de serviço compartilhem a mesma infraestrutura física sob requisitos heterogêneos de vazão, latência, confiabilidade e escala de conectividade. Em um cenário típico, slices eMBB priorizam alta vazão, slices URLLC priorizam latência e confiabilidade, e slices MTC priorizam conectividade massiva e atendimento de tráfego distribuído. Essa heterogeneidade torna a alocação de recursos de rádio um problema de decisão multiobjetivo e dependente do estado da rede.

Arquiteturas O-RAN favorecem o uso de controladores programáveis, xApps e loops de controle para ajustar decisões de rádio de forma automatizada. Nesse contexto, uma estratégia prática consiste em manter o scheduler intra-slice tradicional e otimizar, no nível inter-slice, pesos que controlam a partição do orçamento de recursos entre fatias. Essa separação preserva a compatibilidade com schedulers conhecidos e reduz a complexidade do problema de decisão: em vez de redefinir o scheduler, a busca procura um vetor contínuo de pesos capaz de produzir bom compromisso entre satisfação de SLA, atendimento de carga, entrega de pacotes e utilização de recursos.

Este trabalho avalia meta-heurísticas como mecanismos de busca para esse vetor de pesos. O foco não é propor um novo scheduler independente, mas analisar como diferentes métodos de otimização exploram o espaço contínuo de alocação de slices. As meta-heurísticas avaliadas são PSO, GA, SA e uma abordagem híbrida reformulada. AQPS e RSLAQ são mantidos como bases conceituais e comparativas do estudo: AQPS motiva alocação sensível a QoS em cenários multi-serviço, enquanto RSLAQ fornece a perspectiva de controle orientado a SLA em O-RAN. Os schedulers puros RR, BCQI e PF devem permanecer como baselines de comparação na versão experimental final.

As contribuições deste artigo são:

\begin{enumerate}
  \item Formular a distribuição inter-slice como um problema de otimização contínua sobre pesos de alocação.
  \item Avaliar PSO, GA, SA e uma meta-heurística híbrida usando uma função objetivo comum baseada em KPIs de rede e SLA.
  \item Comparar as melhores soluções encontradas em cenários de baixo tráfego, carga normal e congestionamento.
  \item Identificar limitações estatísticas e experimentais que devem ser corrigidas antes de uma submissão final.
\end{enumerate}

\section{Formulação do Problema}

Considere um conjunto de \(S\) slices. Neste estudo, \(S=3\), correspondendo aos slices eMBB, URLLC e MTC. Uma solução candidata é representada pelo vetor contínuo:

\begin{equation}
\mathbf{w} = [w_{eMBB}, w_{URLLC}, w_{MTC}].
\end{equation}

Cada componente \(w_i\) define o peso de alocação associado ao slice \(i\). A implementação normaliza os pesos para respeitar:

\begin{equation}
w_i \geq 0, \quad \forall i \in \{1,\ldots,S\},
\end{equation}

\begin{equation}
\sum_{i=1}^{S} w_i = 1.
\end{equation}

Assim, as meta-heurísticas operam em um domínio contínuo simples, mas sujeito a normalização. Quando um vetor candidato contém valores negativos ou soma inadequada, a função \texttt{normalize\_weights} da implementação remove valores negativos, completa dimensões ausentes e renormaliza o vetor. Caso a soma seja nula, utiliza-se a alocação uniforme \([1/3, 1/3, 1/3]\).

A qualidade de cada vetor é avaliada por simulação. O problema é tratado como maximização:

\begin{equation}
\max_{\mathbf{w}} F(\mathbf{w}),
\end{equation}

em que \(F(\mathbf{w})\) é o score calculado a partir dos KPIs agregados por slice. A função objetivo implementada combina satisfação média de SLA, pior satisfação de SLA entre slices, atendimento médio da carga oferecida, PDR médio e utilização média do orçamento. Penalidades adicionais são aplicadas para atraso médio URLLC acima de 10 ms, starvation do slice MTC e ocupação de buffer no eMBB. Em forma resumida:

\begin{equation}
F(\mathbf{w}) = 100(0{,}35\bar{s}_{SLA} + 0{,}20s_{SLA}^{min} + 0{,}20\bar{s}_{load} + 0{,}15\bar{p}_{PDR} + 0{,}10\bar{u}) - P_{URLLC} - P_{MTC} - P_{eMBB}.
\end{equation}

Essa função favorece soluções que atendem os três slices, evitando que alta vazão em um slice compense completamente a degradação dos demais. A definição de melhor solução neste artigo segue diretamente o maior valor observado de \(F(\mathbf{w})\).

\section{Métodos Avaliados}

\subsection{PSO}

O PSO mantém uma população de partículas, cada uma representando um vetor de pesos. A posição de cada partícula é atualizada por termos de inércia, melhor posição individual e melhor posição global. Após cada atualização, o vetor é normalizado. Na implementação analisada, o PSO utiliza coeficiente de inércia 0,55 e fatores cognitivo/social 1,30. A busca foi executada com quatro iterações e seis partículas, totalizando 24 avaliações por cenário.

\subsection{GA}

O GA usa uma população de vetores de pesos, seleção por elite, recombinação por crossover e mutação. Em cada geração, os indivíduos são avaliados pelo score de simulação; os melhores compõem a elite e geram a próxima população. Na campanha analisada, o GA também realizou 24 avaliações por cenário, decorrentes de quatro iterações e população de seis indivíduos.

\subsection{SA}

O SA inicia a busca a partir de um vetor uniforme e gera candidatos por mutação com temperatura decrescente. Candidatos piores podem ser aceitos com probabilidade dependente da diferença de score e da temperatura. Nesta execução, o SA realizou quatro avaliações por cenário, uma por iteração. Portanto, seus resultados devem ser interpretados com cautela, pois o orçamento de avaliação foi menor que o de GA, PSO e híbrida.

\subsection{Meta-Heurística Híbrida}

A abordagem híbrida combina elementos populacionais e locais. A população inicial inclui candidatos semeados, incluindo pesos uniformes e configurações estruturadas. A cada iteração, o método avalia partículas, mantém um arquivo de elite, aplica refinamento local por mutação do melhor global e aceita refinamentos por critério inspirado em SA. A próxima população é formada por elites, atualizações inspiradas em PSO, recombinações inspiradas em GA e mutações adaptativas. Na campanha analisada, a híbrida realizou 28 avaliações por cenário, pois cada uma das quatro iterações avaliou a população e um candidato local adicional.

\section{Metodologia Experimental}

Os resultados analisados estão em:

\begin{center}
\path{/home/elioth/Documentos/artigo_jussi/ns-o-ran-gym/results_controlled/heuristics_metaheuristics/20260623_123740/}
\end{center}

A campanha contém resultados de meta-heurísticas para três cenários: \texttt{low\_traffic}, \texttt{normal} e \texttt{congestion}. O bloco de heurísticas contém também \texttt{stressed} e \texttt{insufficient\_resources}, mas esses cenários não possuem resultados meta-heurísticos agregados no mesmo diretório; por isso, não são usados na comparação principal entre meta-heurísticas.

As simulações utilizam um gNB, duplexação TDD, RLC UM, padrão TDD solicitado \texttt{D|D|8D|4GB|4U|U|U}, três slices e scheduler intra-slice PF para os modos slice-aware e \texttt{slice\_custom}. A duração de simulação registrada nos metadados é de 10 s, com 0,2 s de drenagem.

Os cenários comparáveis são apresentados na Tabela~\ref{tab:cenarios}.

\begin{table}[ht]
\centering
\caption{Cenários experimentais comparáveis.}
\label{tab:cenarios}
\begin{tabularx}{\textwidth}{l r Y r r r}
\toprule
Cenário & UEs & UEs por slice & Carga eMBB & Carga URLLC & Carga MTC \\
\midrule
\texttt{low\_traffic} & 10 & eMBB=2, URLLC=2, MTC=6 & 55 Mbps & 1 Mbps & 1 Mbps \\
\texttt{normal} & 20 & eMBB=5, URLLC=5, MTC=10 & 70 Mbps & 1 Mbps & 2 Mbps \\
\texttt{congestion} & 60 & eMBB=15, URLLC=10, MTC=35 & 180 Mbps & 5 Mbps & 60 Mbps \\
\bottomrule
\end{tabularx}
\end{table}

O SLA mínimo configurado é 50 Mbps para eMBB, 1 Mbps para URLLC e 1 Mbps para MTC. A comparação principal usa o score da função objetivo. Métricas de rede por slice, como throughput, atraso, PDR, satisfação de carga e satisfação de SLA, são usadas para interpretar os scores.

\section{Resultados}

\subsection{Comparação entre Meta-Heurísticas}

A Tabela~\ref{tab:metaheuristicas} resume os scores por método e cenário. Os valores de média, mediana e desvio padrão são calculados sobre avaliações de função objetivo dentro de uma única busca, não sobre execuções independentes. Assim, o melhor score é a métrica primária de qualidade de solução; as estatísticas de dispersão indicam apenas o comportamento interno da busca.

\begin{landscape}
\scriptsize
\begin{longtable}{llrrrrrrl}
\caption{Estatísticas dos scores das meta-heurísticas.}\label{tab:metaheuristicas}\\
\toprule
Cenário & Método & Avaliações & Melhor & Média & Mediana & Desvio & Pior & Pesos do melhor candidato (eMBB, URLLC, MTC) \\
\midrule
\endfirsthead
\toprule
Cenário & Método & Avaliações & Melhor & Média & Mediana & Desvio & Pior & Pesos do melhor candidato (eMBB, URLLC, MTC) \\
\midrule
\endhead
\texttt{low\_traffic} & GA & 24 & 56,676 & 48,285 & 56,506 & 19,704 & -23,688 & (0,669; 0,135; 0,195) \\
\texttt{low\_traffic} & PSO & 24 & 56,568 & 52,635 & 56,355 & 9,987 & 20,284 & (0,571; 0,169; 0,259) \\
\texttt{low\_traffic} & SA & 4 & 56,043 & 55,757 & 55,846 & 0,336 & 55,292 & (0,308; 0,197; 0,495) \\
\texttt{low\_traffic} & Híbrida & 28 & \textbf{56,818} & 53,820 & 56,504 & 14,122 & -18,218 & (0,692; 0,174; 0,135) \\
\texttt{normal} & GA & 24 & 67,479 & 56,159 & 61,830 & 11,812 & 37,708 & (0,572; 0,092; 0,336) \\
\texttt{normal} & PSO & 24 & \textbf{68,085} & 57,376 & 63,021 & 19,990 & -25,780 & (0,753; 0,019; 0,228) \\
\texttt{normal} & SA & 4 & 62,108 & 61,182 & 61,296 & 0,885 & 60,029 & (0,278; 0,119; 0,603) \\
\texttt{normal} & Híbrida & 28 & 67,246 & 62,644 & 66,152 & 8,780 & 38,167 & (0,607; 0,127; 0,266) \\
\texttt{congestion} & GA & 24 & 82,991 & 71,898 & 78,156 & 13,051 & 41,468 & (0,526; 0,118; 0,355) \\
\texttt{congestion} & PSO & 24 & \textbf{84,265} & 65,398 & 74,852 & 28,194 & -44,472 & (0,490; 0,062; 0,448) \\
\texttt{congestion} & SA & 4 & 81,869 & 70,362 & 69,427 & 8,718 & 60,723 & (0,399; 0,214; 0,387) \\
\texttt{congestion} & Híbrida & 28 & 81,953 & 75,380 & 77,739 & 9,319 & 43,209 & (0,424; 0,203; 0,373) \\
\bottomrule
\end{longtable}
\end{landscape}

No cenário \texttt{low\_traffic}, a híbrida encontrou o melhor score, mas a diferença para GA e PSO foi pequena: 0,142 ponto acima do GA e 0,250 ponto acima do PSO. Esse cenário mostra uma região de soluções de desempenho muito próximo, na qual pequenas variações de peso alteram pouco o score final.

No cenário \texttt{normal}, o PSO apresentou o melhor score, 68,085, seguido por GA, 67,479, e híbrida, 67,246. A diferença entre os três melhores métodos foi inferior a 0,9 ponto, mas o SA ficou abaixo, com melhor score de 62,108. A melhor solução do PSO deslocou a maior parte do peso para eMBB, mas manteve parcela relevante para MTC, sacrificando URLLC sem violar totalmente sua satisfação de SLA.

No cenário \texttt{congestion}, o PSO apresentou a melhor solução global da campanha, com score de 84,265. O GA alcançou 82,991 e a híbrida 81,953. A superioridade do PSO foi mais clara nesse cenário, em que o vetor encontrado alocou aproximadamente 49,0\% para eMBB, 6,2\% para URLLC e 44,8\% para MTC. Essa configuração reflete a necessidade de sustentar MTC sob carga elevada, sem abandonar completamente eMBB e URLLC.

\subsection{Métricas dos Melhores Candidatos}

A Tabela~\ref{tab:kpis} apresenta os KPIs por slice para os melhores candidatos de cada cenário.

\begin{landscape}
\scriptsize
\begin{longtable}{lllrrrrrr}
\caption{KPIs por slice dos melhores candidatos meta-heurísticos.}\label{tab:kpis}\\
\toprule
Cenário & Método & Slice & Throughput & Atraso médio & PDR & Atendimento da carga & Satisfação de SLA & Utilização \\
\midrule
\endfirsthead
\toprule
Cenário & Método & Slice & Throughput & Atraso médio & PDR & Atendimento da carga & Satisfação de SLA & Utilização \\
\midrule
\endhead
\texttt{low\_traffic} & Híbrida & eMBB & 28,012 Mbps & 20,439 ms & 50,00\% & 50,93\% & 56,02\% & 100,00\% \\
\texttt{low\_traffic} & Híbrida & URLLC & 1,560 Mbps & 4,969 ms & 100,00\% & 100,00\% & 100,00\% & 100,00\% \\
\texttt{low\_traffic} & Híbrida & MTC & 0,682 Mbps & 8,210 ms & 53,33\% & 68,25\% & 68,25\% & 100,00\% \\
\texttt{normal} & PSO & eMBB & 24,388 Mbps & 1112,893 ms & 34,20\% & 34,84\% & 48,78\% & 100,00\% \\
\texttt{normal} & PSO & URLLC & 0,938 Mbps & 14,864 ms & 60,17\% & 93,84\% & 93,84\% & 100,00\% \\
\texttt{normal} & PSO & MTC & 1,843 Mbps & 3,865 ms & 72,04\% & 92,18\% & 100,00\% & 100,00\% \\
\texttt{congestion} & PSO & eMBB & 81,115 Mbps & 1754,333 ms & 44,24\% & 45,06\% & 100,00\% & 100,00\% \\
\texttt{congestion} & PSO & URLLC & 4,479 Mbps & 5,981 ms & 70,00\% & 89,59\% & 100,00\% & 100,00\% \\
\texttt{congestion} & PSO & MTC & 39,878 Mbps & 19,592 ms & 62,95\% & 66,46\% & 100,00\% & 100,00\% \\
\bottomrule
\end{longtable}
\end{landscape}

Os resultados evidenciam que o score não é equivalente a throughput total. No cenário \texttt{normal}, por exemplo, o PSO aceita baixa satisfação do eMBB para preservar URLLC e MTC. Em \texttt{congestion}, o melhor vetor aumenta a participação do MTC, reduzindo risco de starvation e elevando o score mesmo com PDR ainda limitado. Portanto, a função objetivo favorece compromisso entre slices, não maximização isolada de vazão.

\subsection{Comparação com Heurísticas de Alocação no Mesmo Diretório}

Além das meta-heurísticas, a campanha solicitada contém heurísticas de alocação inter-slice: \texttt{slice\_demand\_greedy}, \texttt{slice\_sla\_greedy}, \texttt{slice\_least\_waste}, \texttt{slice\_qos\_mixed} e \texttt{slice\_random\_vine}. Como a busca meta-heurística agregada disponível foi executada apenas com \texttt{seed=1}, a comparação mais direta usa a mesma semente. A Tabela~\ref{tab:meta-vs-heuristica} compara o melhor candidato meta-heurístico com a melhor heurística da semente 1.

\begin{table}[ht]
\centering
\caption{Melhor meta-heurística versus melhor heurística na semente 1.}
\label{tab:meta-vs-heuristica}
\small
\begin{tabularx}{\textwidth}{lYrYrrr}
\toprule
Cenário & Melhor meta-heurística & Score meta & Melhor heurística seed=1 & Score heurística & Ganho absoluto & Ganho relativo \\
\midrule
\texttt{low\_traffic} & Híbrida & 56,818 & \texttt{slice\_demand\_greedy} & 56,654 & +0,164 & +0,3\% \\
\texttt{normal} & PSO & 68,085 & \texttt{slice\_demand\_greedy} & 65,813 & +2,272 & +3,5\% \\
\texttt{congestion} & PSO & 84,265 & \texttt{slice\_random\_vine} & 70,988 & +13,277 & +18,7\% \\
\bottomrule
\end{tabularx}
\end{table}

Os ganhos são quase nulos em \texttt{low\_traffic}, moderados em \texttt{normal} e expressivos em \texttt{congestion}. Isso sugere que a otimização contínua dos pesos é mais útil quando a rede está sob disputa intensa por recursos. Em baixo tráfego, o espaço de soluções contém vários vetores com score semelhante, reduzindo a vantagem prática de uma busca mais sofisticada.

A Tabela~\ref{tab:heuristicas} mostra, em complemento, as estatísticas das heurísticas sobre três sementes. Esses valores não devem ser comparados diretamente com as meta-heurísticas como teste estatístico, pois as meta-heurísticas ainda não foram repetidas nas mesmas três sementes.

\begin{table}[ht]
\centering
\caption{Scores das heurísticas no diretório analisado, agregados sobre três sementes.}
\label{tab:heuristicas}
\small
\begin{tabularx}{\textwidth}{lYrrr}
\toprule
Cenário & Heurística & Melhor & Média & Desvio \\
\midrule
\texttt{low\_traffic} & \texttt{slice\_demand\_greedy} & 99,163 & 71,821 & 23,726 \\
\texttt{low\_traffic} & \texttt{slice\_sla\_greedy} & 99,160 & 71,497 & 23,996 \\
\texttt{low\_traffic} & \texttt{slice\_qos\_mixed} & 99,161 & 71,378 & 24,086 \\
\texttt{low\_traffic} & \texttt{slice\_random\_vine} & 99,001 & 71,503 & 23,893 \\
\texttt{low\_traffic} & \texttt{slice\_least\_waste} & 98,932 & 70,840 & 24,377 \\
\texttt{normal} & \texttt{slice\_sla\_greedy} & 97,997 & 78,656 & 17,049 \\
\texttt{normal} & \texttt{slice\_qos\_mixed} & 96,115 & 77,761 & 16,295 \\
\texttt{normal} & \texttt{slice\_random\_vine} & 97,981 & 76,818 & 19,290 \\
\texttt{normal} & \texttt{slice\_demand\_greedy} & 89,238 & 75,753 & 12,108 \\
\texttt{normal} & \texttt{slice\_least\_waste} & 82,108 & 72,988 & 8,758 \\
\texttt{congestion} & \texttt{slice\_random\_vine} & 70,988 & 68,533 & 3,697 \\
\texttt{congestion} & \texttt{slice\_sla\_greedy} & 54,486 & 53,881 & 0,552 \\
\texttt{congestion} & \texttt{slice\_qos\_mixed} & 54,268 & 52,976 & 1,388 \\
\texttt{congestion} & \texttt{slice\_demand\_greedy} & 52,112 & 51,623 & 0,565 \\
\texttt{congestion} & \texttt{slice\_least\_waste} & 50,854 & 49,977 & 0,828 \\
\bottomrule
\end{tabularx}
\end{table}

A alta variabilidade das heurísticas em \texttt{low\_traffic} e \texttt{normal} mostra que a semente afeta fortemente os resultados. Portanto, a conclusão sobre superioridade das meta-heurísticas deve permanecer restrita: elas superam as heurísticas na comparação pareada com \texttt{seed=1}, mas ainda não há base para afirmar superioridade estatística geral.

\section{Discussão}

Os resultados apontam três conclusões técnicas.

Primeiro, o PSO encontrou as melhores soluções em dois dos três cenários comparáveis. A melhor solução global foi obtida em congestionamento, com score 84,265. Isso sugere que a dinâmica de melhor pessoal e melhor global do PSO foi eficaz na exploração do espaço contínuo de pesos quando o trade-off entre eMBB, URLLC e MTC se tornou mais severo.

Segundo, a meta-heurística híbrida foi competitiva, mas não dominante. Ela venceu no cenário \texttt{low\_traffic}, mas a margem foi muito pequena. Em \texttt{normal} e \texttt{congestion}, ficou abaixo do PSO. Sua média de score foi alta em alguns cenários, mas essa média reflete a trajetória de uma busca com 28 avaliações, não execuções independentes. Assim, não se deve concluir que a híbrida é estatisticamente superior sem novas rodadas.

Terceiro, o ganho da otimização por pesos depende da condição de carga. Em congestionamento, o melhor PSO superou a melhor heurística da semente 1 em 18,7\%. Esse é o resultado mais relevante da campanha, pois indica que a busca contínua consegue encontrar uma partição de recursos mais adequada do que regras manuais quando há forte competição entre slices. Em baixo tráfego, o ganho foi de apenas 0,3\%, o que sugere que heurísticas simples podem ser suficientes quando a pressão por recursos é menor.

Um ponto metodológico importante é que os resultados ainda não fecham a comparação completa exigida para o artigo final. O diretório analisado não contém, no mesmo protocolo de 10 s e com os mesmos artefatos meta-heurísticos, os schedulers puros RR, BCQI e PF, nem comparações completas com AQPS e RSLAQ. Há execuções posteriores no workspace com esses baselines, mas elas usam outro protocolo experimental e não devem ser misturadas sem controle. Para uma versão de submissão, RR, BCQI, PF, AQPS e RSLAQ devem ser reexecutados ou consolidados no mesmo conjunto de cenários, duração, sementes e função objetivo.

\section{Ameaças à Validade}

As principais ameaças à validade são:

\begin{enumerate}
  \item A busca meta-heurística analisada foi registrada para \texttt{seed=1}; não há 30 execuções independentes por método.
  \item O orçamento de avaliação difere entre métodos: GA e PSO têm 24 avaliações, SA tem 4 e a híbrida tem 28.
  \item As médias das meta-heurísticas são médias de avaliações dentro de uma busca, não médias de execuções independentes.
  \item Os cenários \texttt{stressed} e \texttt{insufficient\_resources} não possuem resultados meta-heurísticos agregados no diretório solicitado.
  \item RR, BCQI, PF, AQPS e RSLAQ não estão completos no mesmo protocolo da campanha \texttt{20260623\_123740}.
  \item Não foram aplicados testes de Friedman, Wilcoxon ou Nemenyi, pois os dados disponíveis não constituem amostras pareadas independentes suficientes.
\end{enumerate}

Essas limitações não invalidam a análise exploratória, mas impedem afirmações fortes de significância estatística.

\section{Conclusão}

Este artigo analisou a otimização contínua dos pesos de alocação de slices em redes 5G/O-RAN usando PSO, GA, SA e uma meta-heurística híbrida. Nos resultados disponíveis da campanha \texttt{20260623\_123740}, o PSO foi a meta-heurística que encontrou as melhores soluções de forma mais consistente: venceu nos cenários \texttt{normal} e \texttt{congestion} e obteve o maior score global, 84,265, no cenário congestionado. A meta-heurística híbrida venceu em \texttt{low\_traffic}, com score 56,818, mas por margem pequena em relação ao GA e ao PSO.

Na comparação pareada com heurísticas de alocação usando \texttt{seed=1}, a melhor meta-heurística superou a melhor heurística em todos os três cenários comparáveis. O ganho foi marginal em baixo tráfego, moderado em carga normal e expressivo em congestionamento. Portanto, a principal evidência experimental é que a otimização meta-heurística dos pesos é mais valiosa sob alta competição por recursos.

A conclusão científica deve ser formulada com cautela: PSO é o melhor método segundo o critério de melhor solução encontrada nos dados analisados, enquanto a híbrida é competitiva, mas ainda não comprovadamente superior. Para a versão final do artigo, é necessário executar múltiplas rodadas independentes por método, equalizar o orçamento de avaliações e incluir RR, BCQI, PF, AQPS e RSLAQ no mesmo protocolo experimental.

\section*{Referências a Consolidar}

\begin{thebibliography}{9}

\bibitem{rslaq}
RSLAQ: Robust SLA-driven 6G O-RAN QoS xApp Using Deep Reinforcement Learning. Dados bibliográficos completos a consolidar a partir do documento-fonte.

\bibitem{aqps}
Adaptive QoS-Aware Priority Scheduling Solution for Dynamic Radio Resource Allocation in Multiservice 5G Network Slicing Environments. Dados bibliográficos completos a consolidar a partir do documento-fonte.

\bibitem{oran-slicing-metaheuristics}
Literatura base sobre O-RAN, xApps, 5G network slicing, PSO, GA e Simulated Annealing. Entradas bibliográficas completas devem ser incluídas conforme o veículo de submissão.

\end{thebibliography}

\end{document}
